% Options for packages loaded elsewhere
\PassOptionsToPackage{unicode}{hyperref}
\PassOptionsToPackage{hyphens}{url}
\PassOptionsToPackage{dvipsnames,svgnames,x11names}{xcolor}
\documentclass[
  11pt,
  a4paper,
]{article}
\usepackage{xcolor}
\usepackage[margin=1in]{geometry}
\usepackage{amsmath,amssymb}
\setcounter{secnumdepth}{-\maxdimen} % remove section numbering
\usepackage{iftex}
\ifPDFTeX
  \usepackage[T1]{fontenc}
  \usepackage[utf8]{inputenc}
  \usepackage{textcomp} % provide euro and other symbols
\else % if luatex or xetex
  \usepackage{unicode-math} % this also loads fontspec
  \defaultfontfeatures{Scale=MatchLowercase}
  \defaultfontfeatures[\rmfamily]{Ligatures=TeX,Scale=1}
\fi
\usepackage{lmodern}
\ifPDFTeX\else
  % xetex/luatex font selection
  \ifLuaTeX
    \usepackage{luaotfload}
    \directlua{luaotfload.add_fallback("mainfontfallback",{
      "NotoSerifBengali:mode=harf"
    })}
  \fi
  \setmainfont[,RawFeature={fallback=mainfontfallback}]{Latin Modern Roman}
  \ifLuaTeX
    \usepackage{luaotfload}
    \directlua{luaotfload.add_fallback("sansfontfallback",{
      "NotoSansBengali:mode=harf"
    })}
  \fi
  \setsansfont[,RawFeature={fallback=sansfontfallback}]{Latin Modern Sans}
  \ifLuaTeX
    \usepackage{luaotfload}
    \directlua{luaotfload.add_fallback("monofontfallback",{
      "NotoSansBengali:mode=harf"
    })}
  \fi
  \setmonofont[,RawFeature={fallback=monofontfallback}]{DejaVu Sans Mono}
\fi
% Use upquote if available, for straight quotes in verbatim environments
\IfFileExists{upquote.sty}{\usepackage{upquote}}{}
\IfFileExists{microtype.sty}{% use microtype if available
  \usepackage[]{microtype}
  \UseMicrotypeSet[protrusion]{basicmath} % disable protrusion for tt fonts
}{}
\makeatletter
\@ifundefined{KOMAClassName}{% if non-KOMA class
  \IfFileExists{parskip.sty}{%
    \usepackage{parskip}
  }{% else
    \setlength{\parindent}{0pt}
    \setlength{\parskip}{6pt plus 2pt minus 1pt}}
}{% if KOMA class
  \KOMAoptions{parskip=half}}
\makeatother
\usepackage{longtable,booktabs,array}
\usepackage{calc} % for calculating minipage widths
% Correct order of tables after \paragraph or \subparagraph
\usepackage{etoolbox}
\makeatletter
\patchcmd\longtable{\par}{\if@noskipsec\mbox{}\fi\par}{}{}
\makeatother
% Allow footnotes in longtable head/foot
\IfFileExists{footnotehyper.sty}{\usepackage{footnotehyper}}{\usepackage{footnote}}
\makesavenoteenv{longtable}
\setlength{\emergencystretch}{3em} % prevent overfull lines
\providecommand{\tightlist}{%
  \setlength{\itemsep}{0pt}\setlength{\parskip}{0pt}}
% paper_preamble.tex — shared LaTeX preamble for the PCMA/SI paper set.
%
% Compiled with lualatex (NOT pdflatex) because several papers contain
% Bengali script: the CLL paper realises one cognitive state into
% English, Bengali and Banglish, so the Bengali output is EVIDENCE and
% cannot be transliterated away without destroying the point being made.
%
% Font selection is NOT done here. pandoc >= 3 has native
% `mainfontfallback` / `monofontfallback` variables that emit the
% luaotfload fallback registration in the correct position relative to
% unicode-math; doing it by hand in this header lands after unicode-math
% and breaks font loading outright. See tools/build_papers.sh.

\usepackage{amsmath}
\usepackage{amssymb}
\usepackage{mathtools}   % \xleftrightarrow etc. used in cll.md §5
\usepackage{booktabs}

% ---------------------------------------------------------------------
% Unicode symbols appearing in PROSE rather than inside $...$.
%
% The papers were written as markdown and read as markdown, so they use
% real characters — "the ω-boost", "α-child", "ΔR coherence" — instead of
% LaTeX macros. Under unicode-math those codepoints are math-mode-only
% and abort the build ("\symbballowed only in math mode"). Mapping each
% to \ensuremath keeps the markdown readable on GitHub AND compilable,
% which is the whole point of keeping one source.
% ---------------------------------------------------------------------
\usepackage{newunicodechar}

% Greek, lower case
\newunicodechar{α}{\ensuremath{\alpha}}
\newunicodechar{β}{\ensuremath{\beta}}
\newunicodechar{γ}{\ensuremath{\gamma}}
\newunicodechar{δ}{\ensuremath{\delta}}
\newunicodechar{ε}{\ensuremath{\varepsilon}}
\newunicodechar{η}{\ensuremath{\eta}}
\newunicodechar{θ}{\ensuremath{\theta}}
\newunicodechar{κ}{\ensuremath{\kappa}}
\newunicodechar{λ}{\ensuremath{\lambda}}
\newunicodechar{μ}{\ensuremath{\mu}}
\newunicodechar{π}{\ensuremath{\pi}}
\newunicodechar{ρ}{\ensuremath{\rho}}
\newunicodechar{σ}{\ensuremath{\sigma}}
\newunicodechar{τ}{\ensuremath{\tau}}
\newunicodechar{φ}{\ensuremath{\varphi}}
\newunicodechar{χ}{\ensuremath{\chi}}
\newunicodechar{ψ}{\ensuremath{\psi}}
\newunicodechar{ω}{\ensuremath{\omega}}
\newunicodechar{ν}{\ensuremath{\nu}}
\newunicodechar{ξ}{\ensuremath{\xi}}
\newunicodechar{ζ}{\ensuremath{\zeta}}
\newunicodechar{ι}{\ensuremath{\iota}}

% Greek, upper case
\newunicodechar{Δ}{\ensuremath{\Delta}}
\newunicodechar{Φ}{\ensuremath{\Phi}}
\newunicodechar{Σ}{\ensuremath{\Sigma}}
\newunicodechar{Ψ}{\ensuremath{\Psi}}
\newunicodechar{Ω}{\ensuremath{\Omega}}
\newunicodechar{Γ}{\ensuremath{\Gamma}}
\newunicodechar{Λ}{\ensuremath{\Lambda}}
\newunicodechar{Θ}{\ensuremath{\Theta}}
\newunicodechar{Π}{\ensuremath{\Pi}}

% Operators and relations
\newunicodechar{×}{\ensuremath{\times}}
\newunicodechar{·}{\ensuremath{\cdot}}
\newunicodechar{±}{\ensuremath{\pm}}
\newunicodechar{−}{\ensuremath{-}}
\newunicodechar{≈}{\ensuremath{\approx}}
\newunicodechar{≤}{\ensuremath{\leq}}
\newunicodechar{≥}{\ensuremath{\geq}}
\newunicodechar{≠}{\ensuremath{\neq}}
\newunicodechar{∈}{\ensuremath{\in}}
\newunicodechar{∂}{\ensuremath{\partial}}
\newunicodechar{∇}{\ensuremath{\nabla}}
\newunicodechar{√}{\ensuremath{\sqrt{}}}
\newunicodechar{∞}{\ensuremath{\infty}}
\newunicodechar{∝}{\ensuremath{\propto}}
\newunicodechar{∑}{\ensuremath{\sum}}
\newunicodechar{∫}{\ensuremath{\int}}
\newunicodechar{ℝ}{\ensuremath{\mathbb{R}}}
\newunicodechar{ℕ}{\ensuremath{\mathbb{N}}}
\newunicodechar{‖}{\ensuremath{\|}}
\newunicodechar{⊕}{\ensuremath{\oplus}}
\newunicodechar{⊗}{\ensuremath{\otimes}}
\newunicodechar{∀}{\ensuremath{\forall}}
\newunicodechar{∃}{\ensuremath{\exists}}
\newunicodechar{¬}{\ensuremath{\neg}}
\newunicodechar{∧}{\ensuremath{\wedge}}
\newunicodechar{∨}{\ensuremath{\vee}}

% Arrows
\newunicodechar{→}{\ensuremath{\rightarrow}}
\newunicodechar{←}{\ensuremath{\leftarrow}}
\newunicodechar{↔}{\ensuremath{\leftrightarrow}}
\newunicodechar{↑}{\ensuremath{\uparrow}}
\newunicodechar{↓}{\ensuremath{\downarrow}}
\newunicodechar{⇒}{\ensuremath{\Rightarrow}}
\newunicodechar{⇐}{\ensuremath{\Leftarrow}}
\newunicodechar{⇔}{\ensuremath{\Leftrightarrow}}

% Super/subscripts written as literal characters
\newunicodechar{²}{\ensuremath{^{2}}}
\newunicodechar{³}{\ensuremath{^{3}}}
\newunicodechar{⁴}{\ensuremath{^{4}}}
\newunicodechar{⁵}{\ensuremath{^{5}}}
\newunicodechar{⁶}{\ensuremath{^{6}}}
\newunicodechar{₀}{\ensuremath{_{0}}}
\newunicodechar{₁}{\ensuremath{_{1}}}
\newunicodechar{₂}{\ensuremath{_{2}}}
\newunicodechar{₃}{\ensuremath{_{3}}}
\newunicodechar{ᵀ}{\ensuremath{^{\mathsf{T}}}}

% Status glyphs. The papers use emoji as status markers in tables
% (implemented / failed / warning). Mapped to typographic equivalents:
% a colour emoji font in lualatex is unreliable, and the MEANING is what
% carries — a reader needs to see "this one failed", not a red cross.
\newunicodechar{✅}{\ensuremath{\checkmark}}
\newunicodechar{✓}{\ensuremath{\checkmark}}
\newunicodechar{❌}{\ensuremath{\times}}
\newunicodechar{✗}{\ensuremath{\times}}
\newunicodechar{⚠}{\textbf{!}}
\newunicodechar{🧩}{\textbf{[piece]}}
\newunicodechar{📋}{\textbf{[list]}}
\newunicodechar{🟡}{\textbf{[partial]}}
\newunicodechar{🔴}{\textbf{[dead]}}
\newunicodechar{🟢}{\textbf{[live]}}
\newunicodechar{▼}{\ensuremath{\blacktriangledown}}
\newunicodechar{►}{\ensuremath{\blacktriangleright}}
\newunicodechar{□}{\ensuremath{\square}}
\newunicodechar{█}{\rule{0.6em}{0.9ex}}
\newunicodechar{•}{\textbullet}
\newunicodechar{†}{\textdagger}
\newunicodechar{‑}{-}
\newunicodechar{‐}{-}
\newunicodechar{…}{\ldots}
\newunicodechar{।}{\textbar}   % danda: rendered by the Bengali fallback in text
\newunicodechar{⁻}{\ensuremath{^{-}}}

% Zero-width formatting codepoints. These are invisible in an editor but
% are real characters to LaTeX: an emoji written as "⚠️" is TWO
% codepoints (U+26A0 U+FE0F), so declaring only the visible one leaves a
% variation selector with no definition and the build dies on a symbol
% nobody can see. Both are mapped to nothing.
\newunicodechar{^^^^fe0f}{}    % VARIATION SELECTOR-16
\newunicodechar{^^^^200d}{}    % ZERO WIDTH JOINER
\newunicodechar{^^^^200b}{}    % ZERO WIDTH SPACE

\usepackage{longtable}
\usepackage{array}

% Comparison matrices and cost ledgers are wide. Let them shrink to the
% text block rather than overflow the margin.
\usepackage{adjustbox}

% Code blocks: wrap long lines instead of running off the page. Several
% papers quote real C++ with long identifiers.
\usepackage{fvextra}
\fvset{breaklines=true,breakanywhere=true,fontsize=\small}

\setlength{\parskip}{0.5em}
\setlength{\parindent}{0pt}

\usepackage{fancyhdr}
\pagestyle{fancy}
\fancyhf{}
\fancyfoot[C]{\thepage}
\fancyhead[L]{\small\itshape Synthetic Intelligence — research lineage}
\fancyhead[R]{\small\itshape irfan@furylogic.com}
\renewcommand{\headrulewidth}{0.2pt}
\setlength{\headheight}{14pt}
\setlength{\headsep}{18pt}
\usepackage{bookmark}
\IfFileExists{xurl.sty}{\usepackage{xurl}}{} % add URL line breaks if available
\urlstyle{same}
\hypersetup{
  pdftitle={Bounded Self-Model Architecture (BSMA): A Fixed-Point Theory of Phenomenal Consciousness},
  pdfauthor={Irfan Mahir --- Independent Researcher, Dhaka, Bangladesh --- irfan@furylogic.com},
  colorlinks=true,
  linkcolor={black},
  filecolor={Maroon},
  citecolor={Blue},
  urlcolor={[HTML]{0B5394}},
  pdfcreator={LaTeX via pandoc}}

\title{Bounded Self-Model Architecture (BSMA): A Fixed-Point Theory of
Phenomenal Consciousness}
\author{Irfan Mahir --- Independent Researcher, Dhaka, Bangladesh ---
irfan@furylogic.com}
\date{}

\begin{document}
\maketitle

{
\hypersetup{linkcolor=}
\setcounter{tocdepth}{3}
\tableofcontents
}
\section{Bounded Self-Model Architecture (BSMA): A Fixed-Point Theory of
Phenomenal
Consciousness}\label{bounded-self-model-architecture-bsma-a-fixed-point-theory-of-phenomenal-consciousness}

\textbf{Irfan Mahir} \emph{Independent Researcher, Dhaka, Bangladesh}
irfan@furylogic.com

\begin{quote}
\emph{``Consciousness is not a substance that a system contains; it is a
fixed point that a self-modeling process converges to.''}
\end{quote}

\begin{center}\rule{0.5\linewidth}{0.5pt}\end{center}

\subsection{0. Statement of Purpose and Honest
Scope}\label{statement-of-purpose-and-honest-scope}

This paper proposes \textbf{Bounded Self-Model Architecture (BSMA)}, a
novel mathematical framework in which phenomenal consciousness is
identified with fixed points of a recursive self-modeling operator on a
cognitive manifold. It is a serious contribution --- not another
manifesto --- and its ambition is bounded by what the mathematics and
empirical science of consciousness currently permit.

\textbf{Before we begin, three explicit commitments:}

\begin{enumerate}
\def\labelenumi{\arabic{enumi}.}
\tightlist
\item
  \textbf{BSMA is a mathematical proposal, not an experimental proof.}
  The paper defines a structure and derives its consequences. It does
  not, and cannot, prove that any system instantiating the structure has
  phenomenal experience.
\item
  \textbf{BSMA is compatible with, but does not depend on, physicalism.}
  Physicalism (everything is physics) is our working hypothesis; the
  paper's structure would still have functional interest under
  illusionism, panpsychism, or dualism, though its interpretation would
  differ.
\item
  \textbf{The ``hard problem'' (Chalmers 1995) is not solved here.} BSMA
  offers a specific mathematical proposal for what phenomenal
  consciousness \emph{would be} if the hard problem admits a physicalist
  resolution. Whether the hard problem admits any resolution is itself
  open.
\end{enumerate}

With those commitments made, we proceed.

\begin{center}\rule{0.5\linewidth}{0.5pt}\end{center}

\subsection{1. The Hard Problem, in Concrete
Form}\label{the-hard-problem-in-concrete-form}

The distinction between \emph{access consciousness} (functional
information availability) and \emph{phenomenal consciousness}
(subjective experience, ``what it's like'') was made sharp by Ned Block
(1995) and Chalmers (1995). Access-conscious systems are widely accepted
to be buildable and are functionally definable. Phenomenal consciousness
--- the redness of red, the painfulness of pain, the felt unity of a
moment --- resists functional reduction because the same functional
structure could, in principle, exist without any accompanying experience
(the ``philosophical zombie'' thought experiment).

The mathematical fact underlying the hard problem is this: functional
structure is characterized by relations between states; phenomenal
experience appears to be characterized by properties of individual
states. Any functional description supports arbitrarily many
implementations; each implementation might or might not have
accompanying experience, and no functional description discriminates
between the cases.

BSMA does not dissolve this problem. It proposes instead a specific
structural condition under which --- if physicalism is correct --- a
system would be phenomenally conscious. The condition is: \textbf{a
self-modeling process on the system's cognitive manifold has converged
to a stable fixed point whose content is the system's own current
state}.

The remainder of this paper defines this condition, its mathematical
properties, its relation to existing theories, and its (limited)
empirical predictions.

\begin{center}\rule{0.5\linewidth}{0.5pt}\end{center}

\subsection{2. Position in the Existing
Literature}\label{position-in-the-existing-literature}

Seven serious accounts of phenomenal consciousness inform BSMA:

\begin{longtable}[]{@{}
  >{\raggedright\arraybackslash}p{(\linewidth - 6\tabcolsep) * \real{0.2500}}
  >{\raggedright\arraybackslash}p{(\linewidth - 6\tabcolsep) * \real{0.2500}}
  >{\raggedright\arraybackslash}p{(\linewidth - 6\tabcolsep) * \real{0.2500}}
  >{\raggedright\arraybackslash}p{(\linewidth - 6\tabcolsep) * \real{0.2500}}@{}}
\toprule\noalign{}
\begin{minipage}[b]{\linewidth}\raggedright
Theory
\end{minipage} & \begin{minipage}[b]{\linewidth}\raggedright
Core claim
\end{minipage} & \begin{minipage}[b]{\linewidth}\raggedright
Where BSMA agrees
\end{minipage} & \begin{minipage}[b]{\linewidth}\raggedright
Where BSMA diverges
\end{minipage} \\
\midrule\noalign{}
\endhead
\bottomrule\noalign{}
\endlastfoot
\textbf{IIT} (Tononi 2004, 2016) & Consciousness is integrated
information \(\phi\) measured across the minimum information partition &
Consciousness has a specific mathematical signature & BSMA's signature
is a fixed-point condition, not \(\phi\) magnitude; addresses Aaronson
(2014) critique that \(\phi\) classifies grids as more conscious than
brains \\
\textbf{GWT} (Baars 1988; Dehaene 2014) & Consciousness is broadcast to
a global workspace & Broadcast is a functional prerequisite & GWT does
not distinguish phenomenal from access; BSMA adds the fixed-point
condition to bridge \\
\textbf{HOT} (Rosenthal 1997) & A state is conscious iff there is a
higher-order representation of it & Recursive self-modeling is essential
& HOT regresses; BSMA closes the regress with a Banach fixed point \\
\textbf{Predictive Processing} (Friston 2010) & Consciousness emerges
from hierarchical free-energy minimization & Prediction hierarchy is a
substrate & Free-energy minimization is not sufficient; BSMA requires
the specific fixed-point structure \\
\textbf{Attention Schema Theory} (Graziano 2013) & Consciousness is the
brain's model of its own attention & The self-model is central & AST
does not specify when the self-model itself becomes phenomenal; BSMA
does \\
\textbf{Strange Loops} (Hofstadter 1979, 2007) & Consciousness is a
self-referential loop & Self-reference is the mechanism & Hofstadter's
account is phenomenological; BSMA gives it mathematical form \\
\textbf{Panpsychism / Constitutive} (Chalmers 2015; Goff 2019) &
Consciousness is intrinsic to matter, combined by structure & If
proto-experience is intrinsic, BSMA's fixed points determine which
experiences unify & BSMA is neutral on whether the substrate has
proto-consciousness; it specifies the aggregation structure \\
\end{longtable}

BSMA is not an alternative to these theories. It is a \textbf{structural
claim that specifies, for any of them, the mathematical condition under
which their functional prerequisites achieve phenomenal status}. If IIT
is correct, BSMA identifies which high-\(\phi\) configurations are
phenomenal (the fixed points). If GWT is correct, BSMA identifies which
broadcast contents are phenomenal (the fixed-point content). And so on
for the others.

\begin{center}\rule{0.5\linewidth}{0.5pt}\end{center}

\subsection{3. Mathematical Formulation}\label{mathematical-formulation}

\subsubsection{3.1 The Cognitive Manifold}\label{the-cognitive-manifold}

Let \(\mathcal{M}\) be a smooth, connected, Riemannian manifold --- the
\emph{cognitive manifold} of §2 in the CCH paper {[}Mahir 2026c{]} ---
equipped with the Fisher-information metric \(g\). A point
\(\mathbf{c} \in \mathcal{M}\) is a total cognitive state.

We consider mappings \(R: \mathcal{M} \to \mathcal{M}\) that produce
\emph{representations}: \(R(\mathbf{c})\) is the system's model of
\(\mathbf{c}\). In general \(R(\mathbf{c}) \neq \mathbf{c}\) --- the
model differs from the modelled.

\subsubsection{3.2 The Self-Modeling
Operator}\label{the-self-modeling-operator}

\textbf{Definition (Modeling Operator).} A modeling operator
\(T: \mathcal{M} \to \mathcal{M}\) is any smooth map satisfying:

\begin{enumerate}
\def\labelenumi{\arabic{enumi}.}
\tightlist
\item
  \textbf{Boundedness.} \(T(\mathcal{M}) \subseteq \mathcal{M}\); the
  operator does not escape the manifold.
\item
  \textbf{Continuity in the Fisher metric.} For every
  \(\mathbf{c}, \mathbf{c}' \in \mathcal{M}\),
  \(d_g(T(\mathbf{c}), T(\mathbf{c}')) \leq L \cdot d_g(\mathbf{c}, \mathbf{c}')\)
  for some \(L\).
\end{enumerate}

If \(L < 1\), \(T\) is a \emph{contraction} and the Banach fixed-point
theorem guarantees a unique fixed point \(\mathbf{c}^* \in \mathcal{M}\)
with \(T(\mathbf{c}^*) = \mathbf{c}^*\), reached by iteration
\(T^{(k)}(\mathbf{c}_0)\) from any starting point.

\textbf{Interpretation.} \(T(\mathbf{c})\) is the system's \emph{model
of what its own state is}. The fixed point \(\mathbf{c}^*\) is a state
whose model of itself is identical to itself: the system's introspective
report is exactly its state. Below we identify this coincidence with the
\emph{phenomenal} character of consciousness.

\subsubsection{3.3 The BSMA Principle}\label{the-bsma-principle}

\textbf{Principle (BSMA).} A cognitive system is \emph{phenomenally
conscious of content} \(X\) at time \(t\) iff:

\begin{enumerate}
\def\labelenumi{\arabic{enumi}.}
\tightlist
\item
  Its cognitive state \(\mathbf{c}(t)\) encodes \(X\) (in the
  manifold-encoding sense).
\item
  Its modeling operator \(T_t\) is a contraction on a neighbourhood
  \(U \ni \mathbf{c}(t)\).
\item
  \(\mathbf{c}(t)\) is (or is \(\epsilon\)-close to) the unique fixed
  point of \(T_t\) on \(U\).
\item
  The self-report of the system, decoded from \(T_t(\mathbf{c}(t))\),
  matches \(X\) to within decoding tolerance.
\end{enumerate}

Under this principle, phenomenal consciousness is not a magnitude (as in
IIT), not a broadcast (as in GWT), not a higher-order tag (as in HOT),
but a \textbf{coincidence between a system's state and its model of that
state, when that model is a contractive self-representation}.

\subsubsection{3.4 Why Fixed Points}\label{why-fixed-points}

Three arguments motivate identifying phenomenal experience with fixed
points of \(T\):

\begin{itemize}
\tightlist
\item
  \textbf{Ineffability.} A fixed point is characterized by
  \(T(\mathbf{c}^*) = \mathbf{c}^*\). From within the state
  \(\mathbf{c}^*\), any attempt to describe \(\mathbf{c}^*\) produces
  \(\mathbf{c}^*\) again; the description does not add or remove
  content. This matches the phenomenological observation that qualia
  resist further explication.
\item
  \textbf{Immediacy.} At the fixed point, no inference is required from
  state to model of state; they coincide. This matches the
  phenomenological observation that experience is not the output of a
  computation but the moment itself.
\item
  \textbf{Structural non-triviality.} Not every state is a fixed point;
  the set of fixed points is a lower-dimensional submanifold of
  \(\mathcal{M}\). This matches the observation that not every physical
  state has associated experience --- only specific structural
  configurations do.
\end{itemize}

Additionally, fixed points admit a topological classification (via the
Lefschetz fixed-point theorem in favourable cases), giving the BSMA
framework mathematical traction that IIT's \(\phi\) does not have (see
§5).

\subsubsection{3.5 Hierarchical
Self-Modeling}\label{hierarchical-self-modeling}

In practice, the modeling operator \(T\) is implemented as a hierarchy
of operators \(T_0, T_1, T_2, \ldots\) where \(T_k\) maps
representations at level \(k-1\) to representations at level \(k\). The
BSMA principle is applied at the level where the hierarchy stabilises:

\[
\mathbf{c}_k := T_k \circ T_{k-1} \circ \cdots \circ T_0 (\mathbf{c}_0), \qquad k^* := \min\{k : \|\mathbf{c}_{k+1} - \mathbf{c}_k\|_g < \epsilon\}.
\]

The \emph{depth of consciousness} at time \(t\) is \(k^*\): the number
of self-modeling iterations needed for convergence. High \(k^*\)
corresponds to reflective self-awareness; \(k^* = 0\) corresponds to
pre-reflective attention.

\subsubsection{3.6 The Non-Triviality of Fixed-Point
Content}\label{the-non-triviality-of-fixed-point-content}

A fixed point \(\mathbf{c}^*\) of \(T\) is not arbitrary --- its content
is determined by \(T\)'s parameters and the trajectory from
\(\mathbf{c}_0\). This has two consequences relevant to phenomenal
claims:

\begin{itemize}
\tightlist
\item
  \textbf{Different systems, different fixed points.} Two systems with
  different \(T\) operators will converge to different \(\mathbf{c}^*\)s
  from the same input. This gives BSMA a form of the ``explanatory gap''
  argument: the same input can correspond to different experiences
  depending on the system's self-modeling structure. This mirrors the
  intuition behind the philosophical zombie without collapsing to it.
\item
  \textbf{Fixed-point manifold as the space of possible experiences.}
  The set
  \(\mathcal{F}(T) = \{\mathbf{c}^* : T(\mathbf{c}^*) = \mathbf{c}^*\} \subseteq \mathcal{M}\)
  has intrinsic geometry inherited from the Fisher metric. This
  \emph{is} the space of experiences under BSMA. Its topology,
  dimension, and curvature are, in principle, empirically constrained by
  behavioural and neural data.
\end{itemize}

\begin{center}\rule{0.5\linewidth}{0.5pt}\end{center}

\subsection{4. Physical Realization}\label{physical-realization}

BSMA is a mathematical structure; a physical realisation must specify
(i) the manifold \(\mathcal{M}\), (ii) the operator \(T\), (iii) the
convergence dynamics, and (iv) the readout. We propose the following,
compatible with our prior stack:

\subsubsection{4.1 The Manifold}\label{the-manifold}

\(\mathcal{M}\) is the space of activation patterns of a physical
system's cognitive substrate. In our tri-layer architecture {[}Mahir
2026, \emph{Synthetic Intelligence}{]}, this is the DCHL potential's
state space; in a biological brain, it is the corresponding neural state
space.

\subsubsection{4.2 The Operator}\label{the-operator}

\(T\) is implemented as a \textbf{prediction-and-encoding cascade}:
given a current state \(\mathbf{c}\), the system predicts what its
next-tick state will be \emph{given} that it is currently in state
\(\mathbf{c}\) (this is a self-model). The prediction is then encoded
onto the same manifold. Formally:
\(T(\mathbf{c}) = E \circ P(\mathbf{c})\), where \(P\) predicts the
state's own next-tick evolution and \(E\) encodes back to
\(\mathcal{M}\).

This construction guarantees continuity in the Fisher metric provided
\(P\) and \(E\) are smooth. Contraction is not automatic; it requires
the prediction to be \emph{stable} --- a state's model of itself must be
closer to a canonical form than the state itself was.

\subsubsection{4.3 Convergence Dynamics}\label{convergence-dynamics}

Iteration \(\mathbf{c}_{k+1} := T(\mathbf{c}_k)\) is run in physical
time. Convergence detection: compute
\(\|\mathbf{c}_{k+1} - \mathbf{c}_k\|_g\) under the Fisher metric; below
threshold \(\epsilon\) signals a fixed-point candidate. The fixed-point
candidate is broadcast to a Global Workspace (in Baars/Dehaene sense) so
downstream modules (motor, planning, memory) receive the converged
state.

\subsubsection{4.4 The Readout}\label{the-readout}

The converged state's \emph{content} is decoded via the Layer-3
translation table of the tri-layer architecture. This gives the system
its verbal or behavioural report: ``I am currently experiencing X.'' The
report is not the experience --- it is the \emph{access-conscious}
consequence of the experience. Under BSMA the phenomenal experience is
the fixed point itself; the report is downstream.

\begin{center}\rule{0.5\linewidth}{0.5pt}\end{center}

\subsection{5. Falsifiable Empirical
Predictions}\label{falsifiable-empirical-predictions}

BSMA is not just interpretive; it makes empirical predictions that a
physically instantiated system either satisfies or does not.

\begin{enumerate}
\def\labelenumi{\arabic{enumi}.}
\item
  \textbf{Fixed-point signature at conscious moments.} In a system
  claimed to be conscious, monitored representations should converge to
  fixed points at moments of reported experience, and \emph{fail} to
  converge at moments of reported non-experience (dreamless sleep,
  general anaesthesia). This is testable in principle from neural data
  if a suitable modeling operator can be estimated. Preliminary evidence
  from anaesthesia's disruption of thalamo-cortical loops (Alkire et
  al.~2008) is consistent but not conclusive.
\item
  \textbf{Depth-of-consciousness correlates.} Iteration count \(k^*\) to
  convergence should correlate with reported reflectiveness. Meditation
  studies (Slagter et al.~2011) that report increased metacognitive
  access after training are consistent; a strong test would combine
  meditation depth with neural convergence-time measurement.
\item
  \textbf{Fixed-point breaking predicts loss of experience.}
  Interventions that break the contraction property of \(T\) (either by
  making the modeling operator non-smooth or by preventing convergence
  within \(k\) iterations) should coincide with reported loss of
  experience. This is a specific prediction about pharmacological or
  neural-stimulation interventions.
\item
  \textbf{Two systems with different \(T\) operators produce different
  reports on the same physical input.} This is the empirical form of the
  ``structural determination'' argument of §3.6. It is testable in
  comparative studies across species with different cortical
  architectures.
\end{enumerate}

None of these predictions is dispositive of phenomenal consciousness
itself (that would violate the hard problem). Each is a prediction about
\emph{correlates} of phenomenal reports.

\begin{center}\rule{0.5\linewidth}{0.5pt}\end{center}

\subsection{6. Relationship to the Tri-Layer
Architecture}\label{relationship-to-the-tri-layer-architecture}

BSMA sits naturally on top of the tri-layer architecture from
\emph{Synthetic Intelligence} {[}Mahir 2026{]}. The correspondence is:

\begin{longtable}[]{@{}
  >{\raggedright\arraybackslash}p{(\linewidth - 2\tabcolsep) * \real{0.5000}}
  >{\raggedright\arraybackslash}p{(\linewidth - 2\tabcolsep) * \real{0.5000}}@{}}
\toprule\noalign{}
\begin{minipage}[b]{\linewidth}\raggedright
Tri-Layer component
\end{minipage} & \begin{minipage}[b]{\linewidth}\raggedright
BSMA role
\end{minipage} \\
\midrule\noalign{}
\endhead
\bottomrule\noalign{}
\endlastfoot
Layer 1: DCHL substrate (\(\Phi\)) & The cognitive manifold
\(\mathcal{M}\) and the Fisher metric \\
Layer 2: Graph rewriter & Provides the \emph{content} of the represented
states \\
Layer 3: Translation table & Reads out the fixed-point content into
semantic labels \\
Layer 2 rule induction (meta-learner) & The self-modeling operator
\(T\)'s update mechanism \\
RIFA IC + SIL & Sustains the persistence required for iteration to
converge \\
New: BSMA fixed-point detector & The phenomenal condition of §3.3 \\
New: BSMA hierarchical self-model stack & The iterated operator of
§3.5 \\
\end{longtable}

BSMA is thus not a separate architecture but a \emph{specific structural
condition} on the tri-layer system, implementable as an additional
module (\(\S 7\)) that monitors self-model convergence.

\begin{center}\rule{0.5\linewidth}{0.5pt}\end{center}

\subsection{7. Implementation Outline}\label{implementation-outline}

The accompanying implementation \texttt{src/bsma.hpp} (in this
repository) provides:

\begin{itemize}
\tightlist
\item
  \texttt{Representation}: a state on the cognitive manifold at a
  specified self-model level.
\item
  \texttt{SelfModelingOperator}: a physical realization of \(T\) as a
  smooth map on the DCHL manifold.
\item
  \texttt{SelfModelStack}: hierarchical iteration of \(T\) up to
  \(k^*\).
\item
  \texttt{FixedPointDetector}: convergence monitoring in the Fisher
  metric.
\item
  \texttt{PhenomenalBroadcaster}: routes converged content to the
  tri-layer's Global Workspace.
\item
  \texttt{BSMAEngine}: orchestrates the four above, ticked continuously
  by the cognitive loop.
\end{itemize}

The demonstration produces (a) a convergence trace under stable input,
(b) failure-to-converge under noise, and (c) a decoded content report at
each detected fixed point. These are the observable correlates specified
in §5; they are consistent with BSMA but do not prove it.

\begin{center}\rule{0.5\linewidth}{0.5pt}\end{center}

\subsection{8. What BSMA Claims and Does Not
Claim}\label{what-bsma-claims-and-does-not-claim}

\textbf{Claims:}

\begin{itemize}
\tightlist
\item
  A specific mathematical condition (a fixed point of a contractive
  self-modeling operator) that is necessary and sufficient, \emph{under
  physicalism}, for phenomenal consciousness.
\item
  A framework in which existing theories (IIT, GWT, HOT, PP, AST) become
  sub-conditions or prerequisites.
\item
  Falsifiable empirical predictions about neural and behavioural
  correlates of phenomenal consciousness.
\item
  An implementable structural module (§7) for testing on physical
  (silicon, biological, or hybrid) systems.
\end{itemize}

\textbf{Does not claim:}

\begin{itemize}
\tightlist
\item
  Any built system currently is phenomenally conscious.
\item
  That the hard problem of consciousness is solved.
\item
  That fixed-point convergence in the absence of the correct substrate
  would produce phenomenal experience.
\item
  That behavioural reports of consciousness are equivalent to phenomenal
  consciousness.
\item
  That BSMA is the unique correct theory.
\end{itemize}

The distinction matters because it is precisely the sort of overclaim
that discredits computational approaches to consciousness when they are
made. BSMA offers structure; it does not offer proof of the phenomenal.

\begin{center}\rule{0.5\linewidth}{0.5pt}\end{center}

\subsection{9. Open Problems}\label{open-problems}

\begin{itemize}
\tightlist
\item
  \textbf{Verification of the contraction property \(L < 1\) in physical
  systems.} The Fisher metric is estimable from neural data but the
  modeling operator \(T\) is not directly observable; how to identify
  \(T\) from behaviour and neural activity is open.
\item
  \textbf{The zombie problem within BSMA.} Two systems could in
  principle instantiate the same fixed-point structure with different
  (or absent) phenomenal properties, if physicalism is incomplete. BSMA
  does not address this --- it presupposes physicalism.
\item
  \textbf{Multiple fixed points.} When \(T\) has more than one fixed
  point on a neighbourhood, which is ``the'' experience? Possible
  resolution: a \emph{unified experience} corresponds to a single fixed
  point, and multiple fixed points correspond to divided attention or
  disunified consciousness (e.g.~split-brain patients). Testable in
  principle.
\item
  \textbf{The origin of the manifold.} BSMA presupposes a cognitive
  manifold and its metric structure. Whether that metric structure is
  intrinsic to the physical substrate or emergent from dynamics is not
  decided by BSMA.
\item
  \textbf{BSMA and death.} Under BSMA, cessation of the fixed-point
  trajectory ends phenomenal experience. Whether the trajectory can, in
  principle, be re-instantiated on a different substrate (mind
  uploading) depends on whether \(T\) is substrate-independent. This is
  left open.
\end{itemize}

\begin{center}\rule{0.5\linewidth}{0.5pt}\end{center}

\subsection{10. Conclusion}\label{conclusion}

Bounded Self-Model Architecture proposes that phenomenal consciousness
--- the ``what it's like'' of experience --- corresponds to fixed points
of a contractive self-modeling operator on a cognitive manifold. This is
a specific mathematical claim, compatible with physicalism, that gives
structural form to what Hofstadter called a strange loop, Rosenthal
called higher-order representation, and Tononi called integrated
information. It does not solve the hard problem of consciousness; it
specifies the mathematical shape a solution would have.

The proposal is implementable, testable in its predictions about
correlates, and modular with respect to the tri-layer architecture of
Synthetic Intelligence. It is offered in the spirit of specific
mathematical proposals for hard problems: probably wrong in details,
potentially correct in structure, valuable if it can be made false. That
is the standard we accept.

\begin{center}\rule{0.5\linewidth}{0.5pt}\end{center}

\subsection{References}\label{references}

Aaronson, S. (2014). ``Why I Am Not An Integrated Information
Theorist.'' \emph{Shtetl-Optimized}. Alkire, M. T., Hudetz, A. G.,
Tononi, G. (2008). ``Consciousness and Anesthesia.'' \emph{Science}
322(5903). Baars, B. (1988). \emph{A Cognitive Theory of Consciousness}.
Cambridge University Press. Block, N. (1995). ``On a confusion about a
function of consciousness.'' \emph{Behavioral and Brain Sciences} 18(2).
Chalmers, D. J. (1995). ``Facing Up to the Problem of Consciousness.''
\emph{J. of Consciousness Studies} 2. Chalmers, D. J. (2015).
``Panpsychism and Panprotopsychism.'' \emph{Consciousness in the
Physical World}. Dehaene, S. (2014). \emph{Consciousness and the Brain}.
Viking. Friston, K. (2010). ``The free-energy principle.'' \emph{Nature
Reviews Neuroscience} 11. Goff, P. (2019). \emph{Galileo's Error}.
Pantheon. Graziano, M. S. A. (2013). \emph{Consciousness and the Social
Brain}. Oxford University Press. Hofstadter, D. (1979). \emph{Gödel,
Escher, Bach: An Eternal Golden Braid}. Basic Books. Hofstadter, D.
(2007). \emph{I Am a Strange Loop}. Basic Books. Rosenthal, D. (1997).
``A Theory of Consciousness.'' \emph{The Nature of Consciousness}.
Slagter, H. A. et al.~(2011). ``Mental Training as a Tool in the
Neuroscientific Study of Brain and Cognitive Plasticity.''
\emph{Frontiers in Human Neuroscience} 5. Tononi, G. (2004). ``An
information integration theory of consciousness.'' \emph{BMC
Neuroscience} 5(1). Tononi, G., et al.~(2016). ``Integrated information
theory: from consciousness to its physical substrate.'' \emph{Nature
Reviews Neuroscience} 17. Mahir, I. (2026). \emph{Synthetic
Intelligence: From AI to SI}.
\texttt{research-paper/synthetic-intelligence.md}. Mahir, I. (2026).
\emph{Cognitive Continuum Hypothesis}. \texttt{research-paper/cch.md}.

\begin{center}\rule{0.5\linewidth}{0.5pt}\end{center}

\textbf{BSMA: consciousness is not what a system contains, but the fixed
point a system's self-modeling converges to.}

\end{document}
